\chapter[\emj{🔄}~Recursión y Memoization]{\emj{🔄}~Recursión y Memoization}

\begin{dsaquees}
Las muñecas rusas Matryoshka 🪆. 

Abres una muñeca grande y adentro hay otra igual pero más pequeña. 
Sigues abriendo muñecas hasta que llegas a la más pequeñita de todas,
la cual ya no tiene nada adentro. Ese es tu "punto de parada".

Recursión es una función que se llama a sí misma para resolver una
versión más pequeña del mismo problema. Sigues bajando niveles hasta
llegar al "Caso Base" (la muñeca más chica), y luego vas regresando 
hacia arriba con las respuestas de cada nivel.
\end{dsaquees}
\begin{dsaotraforma}
La Recursión es un paradigma de programación donde una función se
define en términos de sí misma. 

Para que una función recursiva sea correcta, DEBE tener:
1. Caso Base: La condición de salida que detiene las llamadas.
2. Caso Recursivo: La lógica que rompe el problema en una parte más
   pequeña y hace la auto-llamada.

Si olvidas el Caso Base, la función se llamará infinitamente hasta 
agotar la memoria del sistema (Stack Overflow).
\end{dsaotraforma}
\begin{dsadiagrama}
Llamada inicial: fact(3)

[ fact(1) ] -> retorna 1 (¡Caso Base!)
[ fact(2) ] -> espera a fact(1)... recibe 1, retorna 2 * 1 = 2
[ fact(3) ] -> espera a fact(2)... recibe 2, retorna 3 * 2 = 6

Resultado Final: 6
\end{dsadiagrama}
\begin{dsacomofunciona}
- Recursión: Más elegante, fácil de leer para problemas de árboles o 
  gráfos. Usa memoria en el "Call Stack" para cada llamada.
- Iteración (Loops): Más eficiente en memoria (O(1) espacio extra). 
  A veces es más difícil de pensar para problemas complejos.

Memoization (La mejora):
Como vimos en DP, si una función recursiva se llama con los mismos 
argumentos muchas veces, guardamos el resultado en un caché 
(Memoization) para que la segunda llamada sea O(1).
\end{dsacomofunciona}
\section*{PSEUDOCÓDIGO}
\hrulefill
\begin{verbatim}
función factorial(n):
    SI n <= 1: return 1 // Caso Base
    return n * factorial(n - 1) // Caso Recursivo
\end{verbatim}
\section*{CÓDIGO C++}
\hrulefill
\begin{lstlisting}
#include <iostream>
#include <vector>
using namespace std;

// Recursión Simple
long long factorial(int n) {
    if (n <= 1) return 1;
    return n * factorial(n - 1);
}

// Recursión con Memoization (Fibonacci)
long long fibMemo(int n, vector<long long>& memo) {
    if (n <= 1) return n;
    if (memo[n] != -1) return memo[n]; // Revisar caché
    
    return memo[n] = fibMemo(n - 1, memo) + fibMemo(n - 2, memo);
}

int main() {
    cout << "Factorial de 5: " << factorial(5) << endl; // 120
    
    int n = 40;
    vector<long long> memo(n + 1, -1);
    cout << "Fibonacci de " << n << ": " << fibMemo(n, memo) << endl;
    
    return 0;
}
\end{lstlisting}
\begin{dsacomplejidad}
Algoritmo        │ Tiempo (Normal) │ Con Memoization
  ─────────────────┼─────────────────┼────────────────────────
  Factorial        │ O(n)            │ O(n)
  Fibonacci        │ O(2\textasciicircum{}n) 💀       │ O(n) ⚡
  Espacio          │ O(n) (Stack)    │ O(n) (Stack + Cache)
\end{dsacomplejidad}
\begin{dsaentrevistas}

\end{dsaentrevistas}
\begin{dsaresumen}
• Función que se llama a sí misma.
• Requiere Caso Base y Caso Recursivo.
• Úsala para árboles, grafos y problemas de "Divide y Vencerás".
• Memoization: El truco para que la recursión sea ultra rápida.
• Cuidado con el Stack Overflow en inputs grandes.
• Cada nivel de recursión = O(1) de espacio extra en el stack.
════════════════════════════════════════════════════════════════
\end{dsaresumen}

